\documentclass[12pt]{article}
\usepackage[utf8]{inputenc}
\usepackage[T1]{fontenc}
\usepackage{hyperref}
\title{Umot\'ina Language Resources for K-12 Curriculum Development}
\author{Compiled by Claude Sonnet 5 (claude-sonnet-5) for Andr\'as Kornai}
\date{}
\begin{document}
\maketitle
\section{Introduction}
Umot\'ina (Umutina, Barbados; ISO 639-3 umo) is a Bororoan language of western Mato Grosso, Brazil. The community lives in the Terra Ind\'igena Umutina near the Bugres and Paraguay rivers. Sources give 2003 or 2004 as the date of the last fluent speaker's death. In 2015 only two elders were reported to remember the language.
\section{Primary Educational Resources}
A 2015 article in the journal \emph{Ecos} describes the community's efforts to revive the language through a school: \url{https://periodicos.unemat.br/index.php/ecos/article/download/692/668/2290}. I found no textbook or primer.
\section{Linguistic and Reference Materials}
Overview: \url{https://en.wikipedia.org/wiki/Umot%C3%ADna_language}. A numerals page notes the language has a linguolabial consonant, with numerals recorded only from one to five: \url{https://lingweb.eva.mpg.de/channumerals/Umotina.htm}. I found no dictionary or grammar in the search results.
\section{Religious and Community Resources}
A people-group profile is at \url{https://joshuaproject.net/people_groups/15712}.
\section{Digital and Audio Resources}
I found no audio or video resources.
\section{Curriculum Development Context}
Documentation is very limited, so a revitalization curriculum would rest on archival material that still has to be located, plus community knowledge. Searches in Portuguese for ``dicion\'ario Umutina'' would be a sensible next step.
\end{document}